\scanpage{18}
\mspage{- 10 -}\hfill n\handwritten{°358}

sur $K'$. Comme $W'$ engendre l'espace $K$-vectoriel $W$, l'ensemble $H$ des $h\handwritten{\in}R'$
nulles sur $W'$ est aussi l'ensemble des $h\handwritten{\in}R'$ nulles sur $W$ ; comme toute
forme $K'$-linéaire sur $V'$ s'étend de manière unique en un élément de $R'$, il
résulte du th.8 du n\handwritten{°5} du \handwritten{\S\ 4} que $W'$ est l'ensemble des $v\handwritten{\in}V'$ tels que
$h(v)=0$ pour tout $h\handwritten{\in}H$. L'intersection des noyaux des éléments de $H$ est
un sous-espace $K$-vectoriel $W_1$ de $V$ rationnel sur $K'$, et ce qui précède montre que l'on a $W_1\handwritten{\cap}V'=W'=W\handwritten{\cap}V'$, d'o\unsure{\handwritten{\`}{u}|\handwritten{\=}{u}} $W=W_1$ d'après la prop.2. C.Q.F.D.

\uline{Remarques}.- a) On démontre facilement que si $v$ est un vecteur d'un espace
$K$-vectoriel $V$ muni d'une $K'$-structure $V'$, le corps $K'(v)$ est engendré par $K'$
et les scalaires $h(v)$ pour $h$ parcourant l'ensemble des formes linéaires
rationnelles sur $K'$. Soit $W$ un sous-espace $K$-vectoriel de $V$, que nous munirons de $K'$-structure induite par celle de $V$. Comme toute forme linéaire rationnelle sur $K'$ définie sur $W$ se prolonge en une forme linéaire rationnelle
définie sur $V$, on voit que pour tout $v\handwritten{\in}W$, le corps $K'(v)$ calculé ``dans $W$ ''
est égal au corps $K'(v)$ calculé ``dans $V$ ''. Autrement dit, pour tout sous-corps $L$ de $K$ contenant $K'$, on a $W_L=V_L\handwritten{\cap}W$.

b) Lorsque $V$ est de dimension finie, son dual $V^{\handwritten{*}}$ est engendré comme espace
$K$-vectoriel à gauche par l'ensemble $R'$ des formes linéaires rationnelles sur
$K'$. D'après le lemme 2, on voit que $R'$ est alors une $K'$-structure sur $V^{\handwritten{*}}$. On
montre facilement que pour tout sous-corps $L$ de $K$ contenant $K'$, $V_L^{\handwritten{*}}=L.R'$ est
l'ensemble des $f\handwritten{\in}V^{\handwritten{*}}$ tels que $f(V_L)\handwritten{\subset}L$ ; il n'y a donc aucune \unsure{ambiguïté|ambigüité} à
parler dans ce cas d'une ``forme linéaire sur $V$ rationnelle sur $L$ ''. Il n'en
est plus de même lorsque $V$ est de dimension infinie sur $K$ (cf. exerc.\quad ).

6. \uline{Critères de rationalité}.

Si $L$ est un sous-corps de $K$ l'ensemble $K$ considéré comme espace vectoriel
\uline{à gauche} sur $L$ sera noté $ {}_L K$ ; l'ensemble des endomorphismes de cet espace
% archive sheet number: 17
